\section*{Resumen}
\addcontentsline{toc}{section}{Resumen}

Este informe presenta el tercer y último entregable del Trabajo Parcial de Programación Concurrente
y Distribuida y, como pide el enunciado, conserva los dos anteriores con sus correcciones. El caso es
el mismo: K-means sobre los viajes del taxi amarillo de Nueva York \parencite{nyctlc} para
caracterizar patrones de movilidad urbana, en línea con el Objetivo de Desarrollo Sostenible~11.

Los entregables anteriores dejaron un K-means en Go, sin librerías de terceros, en una versión
secuencial y una concurrente con un \emph{worker pool} persistente, acumuladores privados por bloque
y una barrera con \texttt{sync.WaitGroup}; sobre los \cifra{n}~viajes del dataset, la concurrente
alcanza \cifra{speedup-p4}$\times$ con cuatro workers y \cifra{speedup-p8}$\times$ con ocho, con la
misma inercia final para cualquier cantidad de workers.

Este entregable verifica formalmente con Spin la ausencia de deadlocks y la exclusión mutua entre el
coordinador que publica los centroides y los workers que los leen, y agrega una propiedad de progreso
bajo \emph{weak fairness}: la corrida termina. Cada propiedad se expresa como aserción o como fórmula
LTL y tiene un mutante que la viola: un worker que pierde su \texttt{wg.Done} produce el deadlock en
la barrera, y un coordinador que publica sin esperarla viola la exclusión. Una regresión de siete
casos exige, para cada uno, el número y el tipo de error esperados, y corre en la integración
continua. Un análisis del código con Claude Opus 5.5 y un prompt estructurado, contrastado hallazgo
por hallazgo, no encontró errores de sincronización, pero sí dieciocho brechas en el entorno del
algoritmo: un caso de \emph{false sharing} que cuesta un 12\,\% con otros valores de $k$, un harness
de medición sin tests que puede perder una sesión completa, y resultados que no quedan atados al
commit que los produjo.

\textit{Palabras clave:} K-means, programación concurrente, Go, worker pool, Promela, Spin, LTL,
verificación formal, deadlock, exclusión mutua, NYC TLC.
